\section{Inference}
\label{s:inference}

\textbf{The proofs in this section are much more terse and shorter than those in earlier sections. They can be expanded on, or earlier proofs can be contracted.}

\subsection{Test Inversion}
\label{ss:inversion}

Fix a null hypothesis $H_0\colon \beta = \beta_0$ and define
\begin{equation}
  W_i(\beta_0) := Z_i\bigl(Y_i - \beta_0 X_i\bigr), \qquad i = 1,\dots,n.
  \label{eq:Wdef}
\end{equation}
Under $H_0$ we have $W_i(\beta_0) = Z_i\varepsilon_i$, so by \ref{ass:exog} $\E[W_i(\beta_0)] = \E[Z_i\varepsilon_i] = 0$ and by \ref{ass:moments} $\Var(W_i(\beta_0)) = \sigma_{Z\varepsilon}^2 < \infty$.

\citet{dufour1997some} exaplains why test inversion is necessary under weak instruments, since any model where the identification region
\[
  \Theta_0 := \{\theta\in\Theta : \mu_{ZX}(\theta) = 0\}
\]
is non-empty, no confidence set procedure can be simultaneously valid and bounded with probability one.

\subsection{The Median-of-Means Anderson--Rubin Ttest}
\label{sec:mom_ar_test}

We partition the sample into $k$ blocks as in the preceding sections and form, within each block $B_j$, the block mean of the moment~\eqref{eq:Wdef}. Because $W_i(\beta_0) = Z_iY_i - \beta_0 Z_iX_i$ is linear in $\beta_0$, the block mean is
an \emph{affine} function of $\beta_0$:
\begin{equation}
  \bar{W}_j(\beta_0)
    := \frac{1}{m}\sum_{i\in B_j} W_i(\beta_0)
    = \hat{S}_{ZY}^{(j)} - \beta_0\,\hat{S}_{ZX}^{(j)},
    \qquad j = 1,\dots,k.
  \label{eq:Wbar}
\end{equation}
The median-of-means AR statistic is the coordinate-wise median of these block
means,
\begin{equation}
  \widetilde{W}(\beta_0)
    := \med\bigl(\bar{W}_1(\beta_0),\dots,\bar{W}_k(\beta_0)\bigr).
  \label{eq:Wtilde}
\end{equation}


\begin{algorithm}[H]
\label{alg:mom_ar}
\caption{Median-of-Means Anderson--Rubin Test}
\KwIn{Sample $(Y_i,X_i,Z_i)_{i=1}^{n}$, number of blocks $k$, null value
      $\beta_0$, threshold $\tau_n(\delta)$}
\KwOut{Decision: reject or fail to reject $H_0\colon\beta=\beta_0$}
Partition $(Y_i,X_i,Z_i)_{i=1}^{n}$ into $k$ blocks $B_1,\dots,B_k$ of size
$m=\lfloor n/k\rfloor$\;
\For{$j=1,\dots,k$}{
  Compute $\hat{S}_{ZY}^{(j)} = \frac{1}{m}\sum_{i\in B_j} Z_iY_i$ and
          $\hat{S}_{ZX}^{(j)} = \frac{1}{m}\sum_{i\in B_j} Z_iX_i$\;
  Form $\bar{W}_j(\beta_0) = \hat{S}_{ZY}^{(j)} - \beta_0\hat{S}_{ZX}^{(j)}$\;
}
$\widetilde{W}(\beta_0) \leftarrow
  \med\bigl(\bar{W}_1(\beta_0),\dots,\bar{W}_k(\beta_0)\bigr)$\;
\Return{reject if $\abs{\widetilde{W}(\beta_0)} > \tau_n(\delta)$, else fail to
        reject}\;
\end{algorithm}


\begin{theorem}[Size of the MoM-AR test]
\label{thm:mom_ar_size}
Assume \ref{ass:exog}--\ref{ass:moments}. Fix $\delta\in(0,1)$, set
$k=\lceil 8\ln(1/\delta)\rceil$ and $m=\lfloor n/k\rfloor$, and define
\begin{equation}
  \tau_n(\delta) := \sigma_{Z\varepsilon}\sqrt{\frac{32\ln(1/\delta)}{n}}.
  \label{eq:tau}
\end{equation}
Then under $H_0\colon\beta=\beta_0$,
\begin{equation}
  \Prob_{H_0}\!\left[\,\abs{\widetilde{W}(\beta_0)} > \tau_n(\delta)\,\right]
    \le \delta.
  \label{eq:size}
\end{equation}
Consequently the test of Algorithm~\ref{alg:mom_ar} has size at most $\delta$.
\end{theorem}
 
\begin{proof}
Under $H_0$ the random variables $W_i(\beta_0)=Z_i\varepsilon_i$ are i.i.d.\ with mean $\E[W_i(\beta_0)] = 0$ by \ref{ass:exog} and variance $\Var(W_i(\beta_0)) = \sigma_{Z\varepsilon}^2 < \infty$ by \ref{ass:moments}. The statistic $\widetilde{W}(\beta_0)$ is exactly the scalar median-of-means estimator of Algorithm~\ref{alg:mom} applied to the sequence
$(W_i(\beta_0))_{i=1}^{n}$. Theorem~\ref{thm:mom_scalar}, with $\mu=0$ and
$\sigma^2=\sigma_{Z\varepsilon}^2$ and the choice
$k=\lceil 8\ln(1/\delta)\rceil$, gives
\[
  \Prob_{H_0}\!\left[\,\abs{\widetilde{W}(\beta_0) - 0}
       > \sigma_{Z\varepsilon}\sqrt{\tfrac{32\ln(1/\delta)}{n}}\,\right]
       \le \delta,
\]
which is~\eqref{eq:size}.
\end{proof}


\subsection{The MoM-AR Confidence Set}
\label{sec:mom_ar_cs}

Inverting the test of Theorem~\ref{thm:mom_ar_size} produces a confidence set.
\todo{Cite confidence set}

\begin{definition}[MoM-AR confidence set]
\label{def:cs}
For $\delta\in(0,1)$ and $\tau_n(\delta)$ as in~\eqref{eq:tau}, the MoM-AR
confidence set is
\begin{equation}
  CS_{1-\delta}
    := \bigl\{\beta_0\in\R : \abs{\widetilde{W}(\beta_0)}\le\tau_n(\delta)\bigr\}.
  \label{eq:cs}
\end{equation}
\end{definition}


\begin{theorem}[Coverage]
\label{thm:coverage}
Assume \ref{ass:exog}--\ref{ass:moments} and let
$k=\lceil 8\ln(1/\delta)\rceil$. Then
\[
  \Prob_{\beta}\bigl[\beta\in CS_{1-\delta}\bigr] \ge 1-\delta.
\]
\end{theorem}

\begin{proof}
By Definition~\ref{def:cs}, $\beta\in CS_{1-\delta} \iff \abs{\widetilde{W}(\beta)}\le\tau_n(\delta)$. Under the true parameter, $W_i(\beta)=Z_i\varepsilon_i$ has mean zero, so Theorem~\ref{thm:mom_ar_size} (the size bound evaluated at the true null $\beta_0=\beta$) gives $\Prob_\beta\left[\,\abs{\widetilde{W}(\beta)}>\tau_n(\delta)\,\right]\le\delta$. Taking complements yields the result.
\end{proof}

\subsection{Geometry of the Confidence Set}
\label{sec:geometry}

\begin{theorem}[Piecewise-affine structure]
\label{thm:piecewise}
The map $\beta_0\mapsto\widetilde{W}(\beta_0)$ is continuous and piecewise affine. There exist finitely many breakpoints $\beta^{(1)}<\beta^{(2)}<\dots<\beta^{(M)}$, with $M\le\binom{k}{2}$, such that on each open interval of $\R\setminus\{\beta^{(1)},\dots,\beta^{(M)}\}$ the function $\widetilde{W}$ coincides with a single block mean $\bar{W}_j$ and is therefore affine.
\end{theorem}

\begin{proof}
The median of $k$ numbers is the value occupying rank $r:=\lceil k/2\rceil$ once the numbers are sorted. We therefore study how the sorted order of $\bar{W}_1(\beta_0),\dots,\bar{W}_k(\beta_0)$ varies with $\beta_0$.
 
\medskip\noindent\textbf{Step 1 (Crossings are finite).} For a pair $i\neq j$, the equation $\bar{W}_i(\beta_0)=\bar{W}_j(\beta_0)$ reads $\hat{S}_{ZY}^{(i)}-\beta_0\hat{S}_{ZX}^{(i)} =\hat{S}_{ZY}^{(j)}-\beta_0\hat{S}_{ZX}^{(j)}$, i.e.\ $(\hat{S}_{ZX}^{(j)}-\hat{S}_{ZX}^{(i)})\beta_0 =\hat{S}_{ZY}^{(j)}-\hat{S}_{ZY}^{(i)}$. If $\hat{S}_{ZX}^{(i)}\neq\hat{S}_{ZX}^{(j)}$ this has a unique solution; if the slopes are equal the lines are parallel and never cross (or coincide, in which case the pair may be merged). Hence the crossing set $\mathcal{C}:=\{\beta_0 : \bar{W}_i(\beta_0)=\bar{W}_j(\beta_0)\text{ for some } i\neq j\}$ satisfies $\abs{\mathcal{C}}\le\binom{k}{2}$. Enumerate it as $\beta^{(1)}<\dots<\beta^{(M)}$ and set $\beta^{(0)}:=-\infty$, $\beta^{(M+1)}:=+\infty$.
 
\medskip\noindent\textbf{Step 2 (The order is constant between crossings).} Fix an interval $I_\ell:=(\beta^{(\ell)},\beta^{(\ell+1)})$. For any pair $i\neq j$ the difference $D_{ij}(\beta_0):=\bar{W}_i(\beta_0)-\bar{W}_j(\beta_0)$ is affine, hence continuous, and vanishes only on $\mathcal{C}$. Since $I_\ell\cap\mathcal{C}=\varnothing$, $D_{ij}$ has constant sign on $I_\ell$. As this holds for every pair, the total ordering of $\bar{W}_1(\beta_0),\dots,\bar{W}_k(\beta_0)$ is constant on $I_\ell$; let $\pi_\ell$ be the permutation realising it, so that $\bar{W}_{\pi_\ell(1)}(\beta_0)\le\dots\le\bar{W}_{\pi_\ell(k)}(\beta_0)$ throughout $I_\ell$.
 
\medskip\noindent\textbf{Step 3 (The median is one block mean on each interval).} By Step~2 the rank-$r$ value on $I_\ell$ is always supplied by the same index, so $\widetilde{W}(\beta_0)=\bar{W}_{\pi_\ell(r)}(\beta_0)$ for all $\beta_0\in I_\ell$; this is affine. With at most $M+1\le\binom{k}{2}+1$ intervals, $\widetilde{W}$ is piecewise affine.
 
\medskip\noindent\textbf{Step 4 (Continuity).} Let $i=\pi_{\ell-1}(r)$ and $j=\pi_\ell(r)$ be the active median indices on either side of a breakpoint $\beta^{(\ell)}$. If $i=j$ there is nothing to prove. If $i\neq j$, the rank-$r$ slot changes occupant across $\beta^{(\ell)}$, which can occur only if lines $i$ and $j$ exchange ranks there; by continuity of $D_{ij}$ this requires $\bar{W}_i(\beta^{(\ell)})=\bar{W}_j(\beta^{(\ell)})$. Hence the left and right limits of $\widetilde{W}$ at $\beta^{(\ell)}$ coincide with this common value, and $\widetilde{W}$ is continuous. (If three or more lines meet at $\beta^{(\ell)}$, any reshuffling is confined to the indices sharing the common value there, and the same conclusion holds.)
\end{proof}

\begin{corollary}[Finite union of intervals]
\label{cor:union}
$CS_{1-\delta}$ is a finite union of at most $\lceil k/2\rceil + 1$ closed intervals, each possibly unbounded.
\end{corollary}
 
\begin{proof}
The set $CS_{1-\delta}=\{\beta_0:\abs{\widetilde{W}(\beta_0)}\le\tau_n(\delta)\}$ is a sublevel set of the continuous piecewise-affine function $\abs{\widetilde{W}}$, hence a finite union of closed intervals. The componentcount is bounded by the number of times $\widetilde{W}$ enters and leaves the band $[-\tau_n(\delta),\tau_n(\delta)]$, which is at most $\lceil k/2\rceil + 1$ given the at most $\binom{k}{2}+1$ affine pieces and the rank-$r$ selection structure.
\end{proof}

\begin{corollary}[Candidate boundary points]
\label{cor:endpoints}
Every boundary point of $CS_{1-\delta}$ lies in the set
\begin{equation}
  \mathcal{B}_{\mathrm{cand}}
    := \bigcup_{\,j:\,\hat{S}_{ZX}^{(j)}\neq 0}
       \left\{
         \frac{\hat{S}_{ZY}^{(j)}-\tau_n(\delta)}{\hat{S}_{ZX}^{(j)}},\;
         \frac{\hat{S}_{ZY}^{(j)}+\tau_n(\delta)}{\hat{S}_{ZX}^{(j)}}
       \right\},
  \label{eq:candidates}
\end{equation}
a set of at most $2k$ explicitly computable points.
\end{corollary}
 
\begin{proof}
A boundary point satisfies $\abs{\widetilde{W}(\beta_0)}=\tau_n(\delta)$, i.e.\ $\widetilde{W}(\beta_0)=\pm\tau_n(\delta)$. By Theorem~\ref{thm:piecewise} there is an interval on which $\widetilde{W}$ equals one block mean $\bar{W}_j$, so the boundary point solves $\bar{W}_j(\beta_0)=\pm\tau_n(\delta)$ for some $j$. When $\hat{S}_{ZX}^{(j)}\neq 0$, solving $\hat{S}_{ZY}^{(j)}-\beta_0\hat{S}_{ZX}^{(j)}=\pm\tau_n(\delta)$ gives the two values in~\eqref{eq:candidates}. When $\hat{S}_{ZX}^{(j)}=0$ the block mean is the constant $\hat{S}_{ZY}^{(j)}$, which equals $\pm\tau_n(\delta)$ only in the non-generic event $\hat{S}_{ZY}^{(j)}=\pm\tau_n(\delta)$ and so contributes no isolated boundary point.
\end{proof}

\subsection{Monotonicity and single-interval confidence sets}
\label{sec:monotonicity}

The slope of $\widetilde{W}$ on any affine piece is $-\hat{S}_{ZX}^{(j)}$ for the active index $j$. When all block slopes share a sign, $\widetilde{W}$ is monotonic and the confidence set is a single interval.

\begin{proposition}[Deterministic single-interval condition]
\label{prop:mono_det}
If the block sums $\hat{S}_{ZX}^{(1)},\dots,\hat{S}_{ZX}^{(k)}$ all share the same sign, then $\widetilde{W}$ is monotonic on $\R$ and $CS_{1-\delta}$ is a single (possibly empty or unbounded) interval.
\end{proposition}

\begin{proof}
By Theorem~\ref{thm:piecewise} the slope of $\widetilde{W}$ on each piece equals
$-\hat{S}_{ZX}^{(j)}$ for some $j$. If all $\hat{S}_{ZX}^{(j)}>0$ then every slope is non-positive, so the continuous function $\widetilde{W}$ is non-increasing on $\R$; if all $\hat{S}_{ZX}^{(j)}<0$ it is non-decreasing. Since a non-increasing function crosses each of the levels $\pm \tau$ at most once, the band set $\{\beta_0 : -\tau \leq \widetilde{W}(\beta_0) \leq \tau\}$ is convex, and thus a single interval.
\end{proof}

\begin{proposition}[Single-interval condition, Chebyshev]
\label{prop:mono_cheby}
Assume \ref{ass:exog}--\ref{ass:moments} with $\mu_{ZX}\neq 0$. If
\begin{equation}
  m \ge \frac{k\,\sigma_{ZX}^2}{\delta\,\mu_{ZX}^2},
  \label{eq:mono_cheby}
\end{equation}
then with probability at least $1-\delta$ all block sums $\hat{S}_{ZX}^{(j)}$
share the sign of $\mu_{ZX}$; consequently $\widetilde{W}$ is monotone and
$CS_{1-\delta}$ is a single interval.
\end{proposition}
 
\begin{proof}
The sign-consistency hypothesis fails on the event $\mathcal{F}=\bigcup_{j=1}^k\{\hat{S}_{ZX}^{(j)}\le 0\}$. For each block, $\{\hat{S}_{ZX}^{(j)}\le 0\}\subseteq
 \{\abs{\hat{S}_{ZX}^{(j)}-\mu_{ZX}}\ge\mu_{ZX}\}$, and since $\Var(\hat{S}_{ZX}^{(j)})=\sigma_{ZX}^2/m$, Chebyshev's inequality gives $\Prob[\hat{S}_{ZX}^{(j)}\le 0]\le \sigma_{ZX}^2/(m\mu_{ZX}^2)$. The union bound yields $\Prob[\mathcal{F}]\le k\sigma_{ZX}^2/(m\mu_{ZX}^2)$, which is at most $\delta$ precisely when~\eqref{eq:mono_cheby} holds. The conclusion then follows from Proposition~\ref{prop:mono_det}.
\end{proof}
Cantelli can also be used to find the single-interval conditions, but it yields little significant improvement.

\subsection{Computation}
One could compute the confidence set by finding all crossing points between the lines, finding the median line, and then computing if the median line in each interval is between $\pm \tau$.
\todo{write out}